\documentclass[Riemann_Hypothesis.tex]{subfiles}
\begin{document}

\pagebreak
\color{red}
\part{Introduction}
\color{black}

The Riemann hypothesis is one of the most famous and oldest open problems 
in mathematics. Formulated 167 years ago, in 1859, by 
the German mathematician Bernhard Riemann, it remains to this day 
at the heart of research in analytic number theory. 

This problem was introduced in the only paper Riemann published on number 
theory, entitled \textit{Ueber die Anzahl der Primzahlen unter einer gegebenen 
Grösse} (On the Number of Primes Less Than a Given Magnitude). 
In this memoir, Riemann sought to understand the distribution of 
prime numbers and to give an exact expression for the prime counting 
function, often denoted $\pi(x)$.

To do this, Riemann studied the zeta function, initially introduced 
by Leonhard Euler for real variables.

\begin{DefNefEn}[label=df:FonctionZetaEuler_en]{Euler-Riemann zeta function}
    For any complex number $s \in \mathbb{C}$ such that $\Re(s) > 1$, 
    the zeta function is defined by the series:
    \begin{equation*}
        \zeta(s) \;=\; \sum_{n=1}^{+\infty} \frac{1}{n^s}
    \end{equation*}
    Euler had already shown that this series could 
    be expressed as an Eulerian product over the set of 
    prime numbers $\mathbb{P}$:
    \begin{equation*}
        \zeta(s) \;=\; \prod_{p \in \mathbb{P}} \left(1 - p^{-s}\right)^{-1}
    \end{equation*}
\end{DefNefEn}

Riemann's major contribution was to extend the domain of definition of 
this function to the entire complex plane. Through the process of analytic 
continuation, he showed that the $\zeta$ function could be defined for all 
$s \in \mathbb{C} \setminus \{1\}$, where it has a simple pole. 

Moreover, Riemann highlighted a functional equation relating $\zeta(s)$ 
to $\zeta(1-s)$, thus revealing a fundamental symmetry of the function. 
He identified two types of zeros for this function:
\begin{itemize}
    \item The \textbf{trivial zeros}, located at the strictly negative even integers ($-2, -4, -6, \dots$).
    \item The \textbf{non-trivial zeros}, all of which belong to the critical strip defined by $0 < \Re(s) < 1$.
\end{itemize}

The Riemann hypothesis states that all these non-trivial zeros 
have a real part equal to $\frac{1}{2}$.

\begin{theoremEn}[label=th:HypotheseRiemann_en]{Riemann Hypothesis}
    All non-trivial zeros of the $\zeta$ function have a real part of $\frac{1}{2}$. In other words, for any $s \in \mathbb{C}$ such that $0 < \Re(s) < 1$ and $\zeta(s) = 0$, we have:
    \begin{equation*}
        \Re(s) \;=\; \frac{1}{2}
    \end{equation*}
\end{theoremEn}

The importance of this conjecture lies in its intimate connection with prime numbers. 
Riemann indeed proved that the distribution of the non-trivial zeros of the zeta function 
dictates the fluctuations in the distribution of prime numbers.

Following the foundational work of Riemann, the twentieth century saw 
the emergence of new approaches to tackle this problem, notably thanks to the 
major contributions of the French mathematician André Weil. 

In the 1940s, Weil became interested in an analogue of the Riemann hypothesis in 
the context of algebraic geometry over finite fields. He succeeded in proving 
the Riemann hypothesis for algebraic curves over finite fields. These results 
laid the groundwork for the famous Weil conjectures, which were later completely 
resolved by Pierre Deligne.

Besides this breakthrough in algebraic geometry, 
the direct contribution of André Weil 
to classical analytic number theory lies in his generalization of the explicit 
formulas. Riemann had initially established a formula relating prime numbers 
to the zeros of the zeta function. In 1952, according to his 
works (see \cite{weil1952_formules_en}), Weil extended 
this relationship to a broad class of smooth test functions, 
revealing a duality between prime powers 
and the spectrum of non-trivial zeros.

This approach led him to formulate an analytic criterion equivalent 
to the Riemann hypothesis, known today as the Weil 
positivity criterion. To state it, we introduce a specific distribution, 
often called the Weil Hermitian form, acting on an appropriate space of 
test functions.

\begin{DefNefEn}[label=df:FormeHermitienneWeil_intro]{Weil Hermitian form}
    Let $\mathcal{W}$ be the space of Weil test functions on the multiplicative 
    group $\left(\mathbb{R}_+^*, \times\right)$. \par
    Let $f$ and $g$ be two functions in $\mathcal{W}$. \par
    We define $g^*$ as the involutive function defined for any $x \in \mathbb{R}_+^*$ by:
    \begin{equation*}
        g^*(x) \;=\; \frac{1}{x} \, \overline{g\left(\frac{1}{x}\right)}
    \end{equation*}
    We denote by $\Phi = f * g^*$ the multiplicative convolution product of $f$ and $g^*$. \par
    According to André Weil's explicit formulas, we call 
    \textbf{Weil Hermitian form}, which we will denote by $\langle f, g \rangle_{\mathcal{HP}}$, the sesquilinear form defined by the action of the geometric distribution on the test function $\Phi$:
    \begin{equation*}
        \langle f, g \rangle_{\mathcal{HP}} \;:=\; W_{\text{pol}}(\Phi) - W_{\text{arith}}(\Phi) - W_{\text{arch}}(\Phi)
    \end{equation*}
    where the different functionals acting on $\Phi$ are explicitly defined analytically by:
    \begin{align*}
        W_{\text{pol}}(\Phi) \;&=\; \mathcal{M}[\Phi](1) + \mathcal{M}[\Phi](0) \\
        W_{\text{arith}}(\Phi) \;&=\; \sum_{p \in \mathbb{P}} \sum_{n=1}^{\infty} \frac{\ln(p)}{p^{n/2}} \left( \Phi(p^n) + \Phi\left(\frac{1}{p^n}\right) \right) \\
        W_{\text{arch}}(\Phi) \;&=\; \left( \ln(\pi) + \gamma \right) \Phi(1) + \int_{1}^{\infty} \frac{\Phi(x) + \Phi\left(\frac{1}{x}\right) - 2\Phi(1)}{x - \frac{1}{x}} \, \frac{\mathrm{d}x}{x}
    \end{align*}
    with $\mathcal{M}$ denoting the Mellin transform, $\mathbb{P}$ the set of prime numbers and $\gamma$ the Euler-Mascheroni constant.
\end{DefNefEn}

Based on this definition, Weil demonstrated that the validity of the Riemann conjecture 
is equivalent to the positivity of a certain global functional form.

\begin{propositionEn}[label=pr:CriterePositiviteWeil_intro]{Weil positivity criterion}
    Let $\mathcal{Z}$ be the set of non-trivial zeros of
    the Riemann Zeta function.\par
    According to André Weil's theorem~\cite{weil1952_formules_en}, 
    the Hermitian form $\langle \cdot, \cdot \rangle_{\mathcal{HP}}$ 
    has the following properties: \par
    For any function $f \in \mathcal{W}$, the evaluation of this form on the
    diagonal is expressed analytically as a sum over the set
    $\mathcal{Z}$:
    \begin{equation*}
        \langle f, f \rangle_{\mathcal{HP}} \;=\; \sum_{\rho \in \mathcal{Z}} \mathcal{M}[f](\rho) \, \overline{\mathcal{M}[f](1 - \overline{\rho})}
    \end{equation*}
    where $\mathcal{M}[f]$ denotes the Mellin transform of the test function $f$. \par
    Moreover, the Riemann Hypothesis is verified if and only if, for any non-zero 
    function $f \in \mathcal{W}$, we have:
    \begin{equation*}
        \langle f, f \rangle_{\mathcal{HP}} \;\geq\; 0
    \end{equation*}
\end{propositionEn}

This result fundamentally transformed the nature of the problem. 
It allowed translating the Riemann hypothesis into a question of functional 
analysis and spectral theory. This reformulation inspired numerous 
attempts at resolution based on the construction of operators 
whose spectrum would correspond to the zeros of the zeta function, thus 
joining the Hilbert-Pólya spectral conjecture.\medskip\par

The latter suggests that the imaginary parts of the non-trivial zeros of the Riemann 
zeta function correspond to the eigenvalues of a self-adjoint operator 
(a Hamiltonian) acting on a Hilbert space. If such an operator exists, 
the self-adjoint nature guarantees that its eigenvalues are real, which 
immediately implies that all the zeros lie on the critical line, 
thereby proving the Riemann hypothesis.

Subsequently, the French mathematician Alain Connes brought a resolution 
to this program. Relying on noncommutative geometry, 
he constructed an adequate space and established a spectral trace formula 
providing a geometric and spectral realization of the zeros of the zeta function. 
His work thus provided a concrete framework for the Hilbert-Pólya approach.

In the present document, we will review in detail the development of the 
explicit formulas and André Weil's Hermitian form. The main objective 
is to unfold the entirety of the reasoning leading to the proof of the 
Riemann hypothesis in the sought framework. To do this, 
we will show that Weil's positivity criterion 
is indeed satisfied for a specific inner product, of which we will 
give the explicit form.

This approach is articulated in particular around the two results given 
by~\ref{lm:EquivalenceFormeWeil_fr} and~\ref{lm:ConstructionMetriqueHP_fr}, 
which link the Weil distribution to spectral geometry:

\begin{lemmaEn}[label=lm:EquivalenceFormeWeil_intro]{Action of the Weil distribution and spectral equivalence}
    Let $f$ and $g$ be two functions of the Weil space $\mathcal{W}$. \par
    Let the Weil distribution $\mathcal{K}_{\Gamma}$ defined on the product $\mathbb{R}_+^* \times \mathbb{R}_+^*$ by the sum of its components (identity, arithmetic and Archimedean):
    \begin{equation*}
        \mathcal{K}_{\Gamma}(x,y) \;=\; \delta(x-y) \;-\; \sum_{n=1}^{\infty} \frac{\Lambda(n)}{\sqrt{n}} \left[ \delta\left(x - ny\right) + \delta\left(y - nx\right) \right] \;+\; \mathcal{K}_{\infty}(x,y)
    \end{equation*}
    where $\Lambda$ is the von Mangoldt function (Definition~\ref{df:VonMangoldt}), 
    $\mathcal{K}_{\infty}$ the Archimedean kernel (Definition~\ref{df:NoyauArchimedien}), 
    and $\delta$ the Dirac distribution 
    (Definition~\ref{df:DistributionDiracMultiplicative}). \par
    We define the action of this distribution on the tensor product $f \otimes \overline{g}$ with respect to the invariant multiplicative Haar measure $\frac{\mathrm{d}x}{x}$:
    \begin{equation*}
        \langle \mathcal{K}_{\Gamma}, f \otimes \overline{g} \rangle \;:=\; \int_{0}^{\infty} \int_{0}^{\infty} f(x) \overline{g(y)} \, \mathcal{K}_{\Gamma}(x,y) \, \frac{\mathrm{d}x}{x} \frac{\mathrm{d}y}{y}
    \end{equation*}
    According to the work of André Weil~\cite{weil1952_formules_en}, this double integral converges absolutely and satisfies the following spectral equality on the set $\mathcal{Z}$ of non-trivial zeros of the Riemann zeta function:
    \begin{equation*}
        \langle \mathcal{K}_{\Gamma}, f \otimes \overline{g} \rangle \;=\; \sum_{\rho \in \mathcal{Z}} \mathcal{M}[f](\rho) \, \overline{\mathcal{M}[g](1 - \overline{\rho})}
    \end{equation*}
    Consequently, the action of the distribution exactly coincides with the associated
    Weil sesquilinear form $\langle f, g \rangle_{\mathcal{HP}}$,
    described in
    {proposition}~\ref{pr:CriterePositiviteWeil},
    and which was defined in definition~\ref{df:FormeHermitienneWeil}.
\end{lemmaEn}

\begin{lemmaEn}[label=lm:ConstructionMetriqueHP_intro]{Hilbert-Pólya structure and Hamiltonian operator}
    Let $\mathcal{Z}$ be the set of non-trivial zeros of the Riemann zeta function.\par
    According to {lemma}~\ref{lm:EquivalenceFormeWeil},
    the action of the Weil bilinear distribution on the space $\mathcal{W}$
    coincides with the spectral sesquilinear form, which we will
    denote $\langle \cdot, \cdot \rangle_{\mathcal{HP}}$. \par
    According to {proposition}~\ref{pr:CriterePositiviteWeil} (Weil's positivity
    criterion), this Hermitian form satisfies $\langle f, f \rangle_{\mathcal{HP}}
    \geq 0$ for any function $f \in \mathcal{W}$ if and only if the Riemann
    hypothesis is verified. \par
    Under this condition, let $\mathcal{N}$ be the kernel associated with this form
    (isotropic subspace). \par
    Since the form $\langle \cdot, \cdot \rangle_{\mathcal{HP}}$ induces a strictly
    positive definite inner product on the quotient space
    $\mathcal{W} / \mathcal{N}$, we deduce that the topological completion of this
    quotient space generates a separable Hilbert space, which we will denote
    $\mathcal{H}_{HP}$. \par
    On this Hilbert space $\mathcal{H}_{HP}$, there exists an unbounded
    spectral Hamiltonian operator, which we will denote $\mathbf{H}$, whose
    eigenvalues exactly correspond to the imaginary parts of the non-trivial
    zeros of the zeta function. \par
    This operator $\mathbf{H}$ is essentially self-adjoint.
\end{lemmaEn}

The study of this document will conclude with the statement and proof of the following result, 
which will establish the Riemann hypothesis in {theorem}~\ref{th:ResolutionConjectureHP}:

\begin{theoremEn}[label=th:ResolutionConjectureHP_intro]{Geometric spectral realization and proof of the Riemann Hypothesis}
    Let $\mathcal{H}$ be a separable Hilbert space and $\mathbf{D}$ an unbounded, self-adjoint operator with compact resolvent on $\mathcal{H}$. \par
    We assume that the operator $\mathbf{D}$ generates a trace formula satisfying the following spectral absorption identity, for any test function $f \in \mathcal{W}$:
    \begin{equation*}
        \mathrm{Tr}_{\mathrm{reg}}\left( f(\mathbf{D}) f^*(\mathbf{D}) \right) \;=\; \mathcal{K}_\Gamma(f * f^*) \;=\; \sum_{\rho \in \mathcal{Z}} \mathcal{M}[f](\rho) \overline{\mathcal{M}[f](1-\bar{\rho})}
    \end{equation*}
    where $f^*(x) = \overline{f(x^{-1})} x^{-1}$ is the involution on the multiplicative Schwartz algebra, and $\mathrm{Tr}_{\mathrm{reg}}$ denotes the regularized trace on the Hilbert space $\mathcal{H}$. \par
    Then, the Hermitian form $\langle \cdot, \cdot \rangle_{\mathcal{HP}}$ is positive or zero on $\mathcal{W}$, and consequently, the Riemann Hypothesis is true.
\end{theoremEn}

For the hypotheses of the theorem, we will work within a Dixmier 
ideal so that the regularized trace has the properties necessary 
for its realization.


\end{document}